% Ad Atticum 8.6 --- headnote
% ad-atticum-08-06, 21 February 49 BC, Formian villa

\textit{Cicero to Atticus, written from the Formian
villa apparently on the ninth day before the Kalends
of March 49 BC (the manuscript dateline: \textit{Scr.
in Formiano ix K. Mart., ut videtur, a. 705 (49)}).
The 21 February date falls a day or two earlier than
its neighbours in the modern book-numbering (8.5 and
8.7 are both 23 February); the editors have left it
where the manuscripts placed it. Cicero had just
sealed the previous evening's letter --- the one he
meant to send before daybreak --- when the praetor C.
Sosius arrived at Formiae, dropping in on the
neighbouring villa of M'. Lepidus, under whom Sosius
had once served as quaestor. Sosius brought a copy of
a letter from Pompey to the consuls, which Cicero
transcribes in full.}

\textit{Pompey's letter (section 2) reports the
dispatch from L. Domitius at Corfinium and orders the
consuls to bring every available force to a single
rendezvous as fast as they can, leaving only the
minimum garrison at Capua. Section 3 is Cicero's
reaction: \textit{di immortales, qui me horror
perfudit} --- ``Gods immortal, what horror has flooded
through me.'' The clause that follows is corrupt in
the manuscripts (the daggers stand for \textit{nihil
mutasset neglegentia hoc quod cum fortiter et
diligenter tum etiam me hercule}); the sense seems to
be the hope that careful generalship up to now will
not be undone by a sudden lapse. Section 4 is wholly
private --- relief that the quartan fever has left
Atticus, an injunction to Pilia not to keep it longer,
and a small running argument about whether Tiro's
borrowing for expenses is owed to his own scruple or
to Curius's tight-fistedness.}
